\documentclass[11pt,a4paper]{ltjsarticle}
% Common preamble; compile the four documents with LuaLaTeX.
\usepackage{fix-cm}
\usepackage[margin=25mm]{geometry}
\usepackage{luatexja-fontspec}
\setmainjfont{HaranoAjiMincho-Regular.otf}[BoldFont=HaranoAjiMincho-Bold.otf]
\setsansjfont{HaranoAjiGothic-Medium.otf}[BoldFont=HaranoAjiGothic-Bold.otf]
\usepackage{amsmath,amssymb,amsthm,mathtools,booktabs,enumitem}
\usepackage[unicode,hidelinks]{hyperref}
\setlength{\emergencystretch}{2em}
\setlist{nosep}
\newtheorem{theorem}{定理}[section]
\newtheorem{proposition}[theorem]{命題}
\newtheorem{lemma}[theorem]{補題}
\newtheorem{corollary}[theorem]{系}
\theoremstyle{definition}
\newtheorem{definition}[theorem]{定義}
\newtheorem{remark}[theorem]{注意}
\newtheorem{example}[theorem]{例}
\newtheorem{problem}[theorem]{問題}
\newcommand{\C}{\mathbb C}
\newcommand{\Z}{\mathbb Z}
\newcommand{\Q}{\mathbb Q}
\newcommand{\Cat}{\mathcal A}
\newcommand{\DD}{\mathcal D}
\newcommand{\OO}{\mathcal O}
\newcommand{\TT}{\mathcal T}
\newcommand{\PP}{\mathcal P}
\newcommand{\SSer}{\mathbb S}
\newcommand{\Db}{D^{\mathrm b}}
\newcommand{\Hom}{\operatorname{Hom}}
\newcommand{\End}{\operatorname{End}}
\newcommand{\Ext}{\operatorname{Ext}}
\newcommand{\RHom}{\operatorname{RHom}}
\newcommand{\qgr}{\operatorname{qgr}}
\newcommand{\Gr}{\operatorname{Gr}}
\newcommand{\QGr}{\operatorname{QGr}}
\newcommand{\Tor}{\operatorname{Tor}}
\newcommand{\coh}{\operatorname{coh}}
\newcommand{\thick}{\operatorname{thick}}
\newcommand{\gldim}{\operatorname{gldim}}
\newcommand{\GKdim}{\operatorname{GKdim}}
\newcommand{\rk}{\operatorname{rk}}
\newcommand{\id}{\operatorname{id}}
\newcommand{\Hilb}{\operatorname{Hilb}}
\newcommand{\versiondate}{2026年11月17日・v14}

\renewcommand{\versiondate}{2026年11月18日・v18}
\renewcommand{\proofname}{証明}
\newcommand{\R}{\mathbb R}
\newcommand{\HH}{\operatorname{HH}}
\newcommand{\Pic}{\operatorname{Pic}}
\newcommand{\Aut}{\operatorname{Aut}}
\newcommand{\Def}{\operatorname{Def}}
\newcommand{\Perf}{\operatorname{Perf}}
\newcommand{\supp}{\operatorname{supp}}
\hypersetup{pdftitle={Fano圏のSerre根と非可換変形}}
\title{Fano圏のSerre根と非可換変形\\\large Hochschild作用・del Pezzo曲面・安定性・エントロピー}
\author{}
\date{\versiondate\\研究原稿：同時根の形式的持上げ；一般代数化・独立新規性は未確定}
\begin{document}
\maketitle
\begin{abstract}
反Serre関手の根を，数値指数だけでなくHochschild作用と非可換変形によって研究する．
PerryのSerre作用の自明性とTodaの一次変形定理から，指定したexact根の一次持上げ空間はHochschild固定部分となる．
さらにtilting bundleを持ち$\HH^1=\HH^3=0$となる滑らかな射影多様体では，
根付きflat Abel変形が全Artin次数で非障害的であることを示す．
相対Serre関手をtilting代数の双対化双加群から構成し，可逆核の剛性で冪関係を強制する．
相対Serreによる再標識が持上げ集合上で自明になるため，有限巡回平均によって根付き持上げを構成できる．
Fanoの場合には反標準ample pairもArtin変形上に残る．
次数2のGeiser型根を保つ形式変形は可換な幾何学的変形に限られる．
次数4では十個の一次座標部分空間の各方向が根と共に全次数へ延長し，
Fermat三次曲面では可換・非可換方向をともに持つ2次元の形式的な根付き変形問題を得る．
さらに反Serre商が有限な自己同値群の同時持上げを証明する．
次数4の符号二乗根を可換に選べることは，対応するグラフの二部性と同値である．
全五頂点を覆う指定された可換な根の組を保つ形式変形は，可換な幾何学的変形に限られる．
量子平面の基準族とSklyanin族により，三つの三乗根を保つHesse型二次元一次変形空間を明示的に代数化する．
三角形の奇閉路では反Serre商のrank三平行移動格子を記述する．
一般の高次固定部分の同時線型化，次数4・三次曲面の族の代数化，奇閉路の全Artin同時持上げは結論しない．
射影平面の同じ数値根を区別するHochschild指標，Fano全圏でのGepner型非存在，およびエントロピーの計算も与える．
root pairの巡回coherenceと一般非可換Fano認識は別問題として残す．
\end{abstract}
{\small\tableofcontents}

\section{問題と先行研究}
\label{sec:intro}
$\C$上で考える．滑らかな射影$n$次元Fano多様体$X$の標準heartについて
\[
 \tau=\SSer^{-1}[n]=(-)\otimes\omega_X^{-1}
\]
と置く．本稿の根は，$\coh X$上のexact自己同値$s$と自然同型$s^r\simeq\tau$である．
根の同型の選択は，特に断らない限り標識しない．
Bondal--Orlov\cite[Theorem 3.1]{BO}から$s=L\otimes f^*$と書け，
\begin{equation}\label{eq:norm}
 f^r=1,\qquad \bigotimes_{j=0}^{r-1}(f^j)^*L\simeq\omega_X^{-1}
\end{equation}
が根の必要十分条件となる．これは$L^r=\omega_X^{-1}$より広いnorm条件である．
同研究の関連原稿\cite{Index}では，数値作用のunipotenceによる正規化，曲面型擬格子の根，
del Pezzo曲面のexact指数を扱った．本稿では\emph{指定した根が非可換方向へ続くか}を問う．

Kuznetsov--Smirnov\cite[Definition 2.1, Theorem 2.8]{KS}の圏論的偏極は，既に
$s^r\simeq\SSer^{-1}[a]$を用い，残余圏のmutation回転へ引き継がれる．
本研究は根という概念の新設ではない．Minamoto\cite[Definition 2.6, Theorem 2.8]{Minamoto}の
ampleな両側tilting複体，Mori--Ueyama\cite[Theorem 4.1]{MU}のfull geometric helixによる認識，
Loginov--Orlov\cite[Definitions 2.1, 2.6]{LO}の非可換Fano/minifoldは，
heart，消滅，生成の条件が互いに異なる．数値根だけからそれらの仮定を得たことにはしない．
Hanihara\cite[Definition 2.5, Theorem 3.3]{Hanihara}のroot pairには生成条件があり，
Calabi--Yau completionにはさらに巡回不変性を要する．

本稿の一次変形に関する直接の入力は次の二つである．Perry\cite[Proposition 1.2, \S5]{Perry}は
properなenhanced圏のSerre関手が$\HH^*$に自明に作用することを示す．
Toda\cite[Theorem 1.1]{Toda}はFourier--Mukai同値を，そのHochschild作用に対応する一次Abel変形間へ持ち上げる．
Abel圏の変形と$\HH^2,\HH^3$の役割はLowen--Van den Bergh\cite[Theorem 3.1]{LVdB}による．
これら既知定理の組合せを本稿の新規定理と混同しない．
本稿が直接証明・計算する部分は，heart-exact性の継承，冪関係の補正，
del Pezzo曲面上の自然な反標準表現と持上げ部分空間，および安定性・エントロピーとの比較である．
それらの独立した文献上の新規性は未確定である．

全Artin次数への拡張では，Hille--Perling\cite[Theorem 1.1]{HP}の合理曲面上のtilting bundle，
Barmeier--Wang\cite[Theorem 3.18, Proposition 3.21]{BW}のAbel変形とtilting代数変形の比較，
Lowen\cite[Theorems A, B]{LowenOb}の対象・核の障害理論を入力とする．
これら既知定理は，それだけで指定した根の冪関係を与えるものではない．
本稿では，相対Serreの全変形上での構成，$\HH^1=0$による核の一意性，
小拡大の持上げtorsor上での有限次数再標識を組み合わせて，
定理\ref{thm:formal-roots}を証明する．
次数2の形式的可換性（系\ref{cor:geiser-formal}），次数4の形式的存在
（系\ref{cor:degree4-formal}），Fermat三次曲面の2次元変形
（命題\ref{prop:fermat-formal}）がその具体的適用である．
既存の変形理論の適用である部分と，これらの組合せ・明示計算とを区別する．

複数根については定理\ref{thm:finite-package}が有限反Serre商を持つ群の同時変形を扱う．
定理\ref{thm:bipartite-roots}は次数4のPicard格子の整性から二部グラフ判定を与え，
系\ref{cor:graph-formal}が全形式次数の同時変形に適用する．
節\ref{sec:quantum-benchmark}は既知のskew polynomial環とNakayama公式\cite[Example 5.5]{RRZ}を使う基準例であり，
その代数族自体の新規性は主張しない．
節\ref{sec:hesse-sklyanin}ではSklyanin族のKodaira--Spencer写像を計算し，Hesse不変空間の二方向を実現する．
節\ref{sec:triangle-group}では奇閉路のaffine Picard表現を完全に計算する．

Hochschild--Serre代数を用いるLin\cite{Lin}，Lin--Rampazzo--Zhang\cite{LRZ}の研究は，
通常の$\HH^*$に留まらずSerre捩れ係数付き自然変換を使って残余圏から幾何を復元する．
通常の$\HH^*$上でSerre作用が自明であることを，全Hochschild--Serre代数上の作用の自明性と取り違えない．
本稿の固定部分は\emph{変形の対称性}を測り，それ自体はTorelli定理ではない．

\section{Hochschild作用と一次変形}
\label{sec:hh}
\begin{proposition}[有限位数のHochschild作用]\label{prop:finite-hh}
$\mathcal D$をproperな$\C$線形enhanced圏，$s$をenhanced自己同値とし，
$s^r\simeq\SSer^{-1}[n]$とする．$\alpha=s_{\HH}$は$\HH^*(\mathcal D)$上で$\alpha^r=1$を満たす．
特に有限次元の$\HH^2$上ではsemisimpleであり，
\[
 P_s=\frac1r\sum_{j=0}^{r-1}\alpha^j
\]
は不変部分への射影である．
\end{proposition}
\begin{proof}
Serre作用はPerryの定理により恒等である．shiftによる共役もHochschild\emph{コホモロジー}上恒等である：
対角双加群をshiftしてから逆shiftする共役において，evaluationの符号を含めると
対角の自己Extの作用は恒等となる．これはshiftがHochschild\emph{ホモロジー}に作用する際の符号とは異なる．
従って$\alpha^r=1$．標数零では$x^r-1$が重根を持たず，平均化の主張が従う．
\end{proof}

\begin{remark}\label{rem:not-cyclic}
ここで得られるのは$\HH^*$上の$C_r$表現であって，圏上の$C_r$作用ではない．
通常$s^r=\tau\ne\id$である．Perryのequivariantization公式やHaniharaの巡回不変性を
この等式だけから適用しない．また数値$K$群上のunipotenceと$\HH^2$上の作用の自明性は別条件である．
\end{remark}

$R=\C[\varepsilon]/(\varepsilon^2)$とする．滑らかな射影$X$に対し，
$u\in\HH^2(X)$で指定される標識付きの一次Abel変形を$\mathcal A_u$と書く．
これはTodaの$\coh(X,u)$であり，$\varepsilon$で消える対象の部分圏を$\coh X$と同定する．
「持上げ」は$R$線形で，この特殊化と互換なものを意味する．
ここで扱うのは標識付き一次変形であり，自己同値群で割ったmoduliの接空間ではない．

\begin{lemma}[特殊ファイバーのheart-exact性]\label{lem:heart}
$\mathcal B,\mathcal B'$を$R$線形Abel圏とし，$\varepsilon$で消える対象をそれぞれ
$\mathcal A,\mathcal A'$とする．導来同値$\Phi_R:\Db(\mathcal B)\to\Db(\mathcal B')$が
特殊ファイバーのpushforwardと互換で，その制限$\Phi:\Db(\mathcal A)\to\Db(\mathcal A')$がheart-exactなら，
$\Phi_R$もheart-exactである．
\end{lemma}
\begin{proof}
任意の$M\in\mathcal B$には$0\to\varepsilon M\to M\to M/\varepsilon M\to0$があり，
両端は$\varepsilon$で消える．その像は仮定により$\mathcal B'$に属するから，
heartの拡大閉性により$\Phi_R(M)\in\mathcal B'$である．逆同値にも同じ議論を使う．
対象$M$の$R$平坦性は要求していない．
\end{proof}

\begin{theorem}[指定した関手の一次持上げ]\label{thm:fixed}
$X$を滑らかな射影多様体，$s$を$\coh X$上のexact自己同値とする．
$s$が$\mathcal A_u$上のexact自己同値へ持ち上がることと
\begin{equation}\label{eq:fixed}
 s_{\HH}(u)=u
\end{equation}
は同値である．$s$が反Serre根であるとき，この空間は命題\ref{prop:finite-hh}の$\operatorname{im}P_s$である．
ここでの結論は\emph{関手の持上げ}であり，選んだ持上げの$r$乗が指定した相対反Serre関手と同型になることはまだ含まない．
\end{theorem}
\begin{proof}
標識を$s$で変えると変形類は$s_{\HH}u$となる．標識付き一次Abel変形の分類
\cite[\S3]{LVdB}により，持上げの存在には\eqref{eq:fixed}が必要である．
逆にToda\cite[Theorem 1.1]{Toda}は$\Db(\mathcal A_u)$と$\Db(\mathcal A_{s_{\HH}u})$の
Fourier--Mukai同値を，特殊ファイバーのpushforwardと互換に構成する．
\eqref{eq:fixed}なら同じ変形上の同値であり，補題\ref{lem:heart}でexact性を得る．
\end{proof}

\begin{lemma}[冪関係のsquare-zero補正]\label{lem:repair}
群の準同型$G_R\to G$の核$V$が$\C$ベクトル空間の加法群であるとする．
$\tau_R\in G_R$は中心元で，その像を$\tau$とする．$s^r=\tau$と，$s$の持上げ$a\in G_R$があるなら，
ある$b\in V$について$(ba)^r=\tau_R$となる．
\end{lemma}
\begin{proof}
$c=a^r\tau_R^{-1}\in V$と置く．$A=\operatorname{Ad}(a)|_V$とすれば
$Ac=c$である．$V$は可換で$\tau_R$は中心だから$A^r=1$でもある．
加法記法で$b=-c/r$と取ると
\[
 (ba)^r\tau_R^{-1}=(1+A+\cdots+A^{r-1})b+c=0.
\]
よって主張が従う．これは群の同型類レベルの証明である．
\end{proof}

\begin{corollary}[相対Serre仮定付きの根の持上げ]\label{cor:relative-root}
定理\ref{thm:fixed}の設定で，exact自己同値の特殊化群の核が加法的ベクトル群であり，
中心的なexact相対反Serre持上げ$\tau_R$が指定されていると仮定する．
このとき$s^r=\tau$を満たす根が，$s_R^r\simeq\tau_R$を満たす根へ持ち上がるための必要十分条件は\eqref{eq:fixed}である．
\end{corollary}
\begin{proof}
必要性は関手の持上げに対する必要性．十分性は定理\ref{thm:fixed}と補題\ref{lem:repair}による．
\end{proof}

\begin{remark}[条件の意味と限界]\label{rem:relative}
通常のMorita/Fourier--Mukaiのsquare-zero変形問題では恒等の持上げの接空間は$\HH^1$である．
ただし一般の六公理の候補heartでは，任意のAbel変形にexactな相対Serre関手が存在することや，
その自己同値群の核が無条件でこの標準形を持つことは仮定なしには分からない．
系\ref{cor:relative-root}ではこれらを明示的な仮定とする．
$\HH^1=0$で，持上げが標準のperfect kernel変形理論に従う場合，
同じ特殊化を持つ核の一意性から冪関係の確認は特に簡単になる．
次節ではtiltingモデルと$\HH^1=\HH^3=0$の下で全Artin次数の持上げを証明する．
自然同型を含む巡回coherenceはなお別問題である．
$\Db(\mathcal A_u)$を$\C$上のsmooth proper圏とみなしてはならない．
\end{remark}


\section{Tiltingモデルによる全Artin次数の根の持上げ}
\label{sec:artin}
本節では局所Artin $\C$代数$R$，極大ideal $\mathfrak m_R$，剰余体$\C$を用いる．
変形は特殊ファイバーの同定を持つflat Abel変形とし，その同定と互換なexact同値で比較する．
自己同値と自然同型を含む2-groupoidを用い，以下のtorsorと固定点の主張はその同型類について述べる．
根の次数$r$は正整数で，$\C$では可逆である．

\subsection{相対Serre関手の存在とheart}
\begin{lemma}[有限代数モデルの相対双対性]\label{lem:relative-algebra}
$A$を有限次元で有限global dimensionの$\C$代数，$A_R$をその有限自由$R$代数変形とする．
$A_R$は$R$上smooth properであり，
\[
 D_RA_R=\operatorname{Hom}_R(A_R,R)
\]
は両側tilting複体である．$\Perf(A_R)$上の相対Serre関手は
$\mathbb S_R=(-)\otimes^{\mathbf L}_{A_R}D_RA_R$で与えられ，Artin基底変換と互換である．
この関手と逆関手は$D^b(\operatorname{mod}A_R)$にも作用する．
\end{lemma}
\begin{proof}
まず，$R$自由な有限代数$B_R$と有限$R$自由$B_R$加群$M_R$について，
$M_0$が$B_0$上projectiveなら$M_R$もprojectiveである．
$M_0$を有限自由加群の直和因子として書き，そのidempotentをnilpotent idealを越えて持ち上げる．
得られたprojective加群$Q_R$から$M_R$への，特殊化で同型となる射をprojectivityで持ち上げる．
余核はNakayamaで零である．核の特殊化も$M_R$の$R$平坦性から零なので，再びNakayamaで核も零となる．

$B_R=A_R^e=A_R^{\mathrm{op}}\otimes_R A_R$，$M_R=A_R$とする．
有限自由$B_R$加群からの全射を反復して分解を作ると，syzygyは全て$R$自由である．
特殊ファイバー$A$は$A^e$上perfectなので，有限段階のsyzygyの特殊化はprojectiveとなる．
前段落によりそのsyzygy自身がprojectiveであり，対角は$A_R^e$上perfectとなる．
ここで$A^e$上のperfect性は，完全体上の有限次元有限global dimension代数の標準的なsmoothnessである．
有限$R$自由性がpropernessを与える．

同じsyzygyの議論を$D_RA_R$の左右加群へ適用すると，両側でperfectとなる．
そのtensorと$\mathbf R\operatorname{Hom}$からなる随伴のunit・counitは，
特殊ファイバーで有限global dimension代数のSerre同値のunit・counitとなる．
それらのconeは$R$-perfectで特殊化が零なので，derived Nakayamaにより零である．
従って$D_RA_R$は可逆である．有限projective分解上のtensor--Hom随伴から，
$M,N\in\Perf(A_R)$について
\[
 \mathbf R\operatorname{Hom}_R(\mathbf R\operatorname{Hom}_{A_R}(M,N),R)
 \simeq\mathbf R\operatorname{Hom}_{A_R}(N,M\otimes^{\mathbf L}_{A_R}D_RA_R)
\]
を得る．有限分解と$R$自由性により基底変換とも互換する．
核が左右perfectなので，tensorと逆tensorは有界有限生成cohomologyを保つ．
\end{proof}

\begin{lemma}[Artin版のheart判定]\label{lem:artin-heart}
補題\ref{lem:heart}は任意の局所Artin基底$R$について成立する．
\end{lemma}
\begin{proof}
$M$の有限filtration $\mathfrak m_R^iM$を用いる．各subquotientは特殊ファイバーからのpushforwardである．
特殊ファイバーでheart-exactな関手はそれらをheartへ写す．有限回の拡大により$M$もheartへ写り，逆関手も同様である．
\end{proof}

\begin{proposition}[Tilting付き幾何の相対反Serre関手]\label{prop:tilting-serre}
$X$を滑らかな射影$n$次元多様体でtilting bundleを持つものとする．
$\coh X$の任意のflat Artin Abel変形$\mathcal A_R$は，
exact自己同値$\tau_R$を持つ．その導来核は相対Serre関手の逆の$[n]$シフトであり，
特殊化は$(-)\otimes\omega_X^{-1}$である．$\tau_R$は基底変換と互換で，
exact自己同値の同型類群の中心に属する．
\end{proposition}
\begin{proof}
Tilting bundle $E$は$\Ext^{>0}(E,E)=0$によりArtin次数ごとに持ち上がる．
Lowen\cite[Theorems A, B]{LowenOb}の対象・複体の持上げと，
Barmeier--Wang\cite[Theorems 3.1, 3.18, Proposition 3.21]{BW}の変形の比較を用いる．
これにより$\mathcal A_R$の導来圏は，有限自由変形$A_R$ of $A=\operatorname{End}(E)$の
加群導来圏と，特殊化に互換に同値となる．coherent有界部分についても同値である．
実際，特殊ファイバーではtilting同値が有界coherent部分を保ち，
有限$\mathfrak m_R$-filtrationとその逆の議論でこの性質が持ち上がる．
$A$は有限次元有限global dimensionである．

補題\ref{lem:relative-algebra}の核を移送して$\mathbb S_R$を得て，
$\tau_R=\mathbb S_R^{-1}[n]$と置く．特殊ファイバーは通常の反標準tensorなので，
補題\ref{lem:artin-heart}により$\tau_R$はheart-exactである．
相対Serre関手の一意性から，任意の$R$線形自己同値との可換性とtilting表示からの独立性が従う．
\end{proof}

\begin{remark}
$\mathbb S_R$の双対性は\emph{$R$に関する相対perfect部分}での双対性である．
非reducedな$R$上の$D^b(\mathcal A_R)$を$\C$上smooth properとしたり，
全対象に絶対Serre双対性や有限global dimensionを課したりはしていない．
本節の変形比較は引用したtiltingモデルの範囲で行い，一般の六公理の候補heartへ無条件に適用しない．
\end{remark}

\subsection{核の剛性と冪関係}
\begin{lemma}[同一変形上の可逆核の一意性]\label{lem:kernel-rigid}
上の$A$について$\HH^1(A)=0$とする．固定した$A_R$上の二つの可逆核が，
特殊ファイバーで同じ可逆核$P$に同型なら，両者は同型である．
さらに$\HH^0(A)=\C$なら，任意の可逆核の自己同型群は$R^*$である．
\end{lemma}
\begin{proof}
可逆性により$\operatorname{Ext}_{A^e}^j(P,P)\simeq\HH^j(A)$である．
小拡大$R'\twoheadrightarrow R$，$I=\ker(R'\to R)$，$\mathfrak m_{R'}I=0$を取る．
固定した核の二つの持上げの差は$\HH^1(A)\otimes I$のtorsorに属する
\cite[Theorem B]{LowenOb}．この空間が零なので，既に存在する持上げは一意である．
$R$の長さに関する帰納法で第一の主張を得る．同型自体の一意性は主張していない．

$\HH^0(A)=\C$なら$Z(A_R)=R$である．実際，$A_R=R1\oplus V_R$と自由加群として分け，
$V_R$から$\operatorname{End}_R(A_R)$へのcommutator写像を考える．
その特殊化は単射なので，ある最大minorが可逆であり，元の写像も単射である．
従って中心はscalarだけである．可逆核の自己準同型環は中心に同型なので主張を得る．
\end{proof}

\begin{theorem}[関手の持上げは根の持上げになる]\label{thm:artin-power}
命題\ref{prop:tilting-serre}の$X$に$\HH^1(X)=0$を仮定する．
heart-exactな根$s^r=\tau$が$\mathcal A_R$のexact自己同値$s_R$へ持ち上がれば，
\[
 s_R^r\simeq\tau_R
\]
となる．従ってdual numbers上では，根そのものの持上げの必要十分条件は$s_{\HH}u=u$である．
選んだ根は一つの変形上で同型まで一意である．
\end{theorem}
\begin{proof}
Tiltingモデルで$s_R^r$の核と$\tau_R$の核は同じ特殊化を持つ．
両者は可逆なので，補題\ref{lem:kernel-rigid}により同型である．
持上げ$s_R$自身の一意性も同じ補題による．一次の主張は定理\ref{thm:fixed}と合わせる．
\end{proof}

\subsection{有限次数の再標識と全次数の存在}
\begin{theorem}[根を持つArtin変形の形式的滑らかさ]\label{thm:formal-roots}
$X$を連結な滑らかな射影多様体でtilting bundleを持つものとし，
\[
 \HH^1(X)=\HH^3(X)=0
\]
を仮定する．$s^r=(-)\otimes\omega_X^{-1}$をheart-exactな根とする．
$s$を相対根として持ち上げた標識付きflat Abel変形は，小拡大に関して非障害的である．
その一次変形空間は$\HH^2(X)^{s_{\HH}}$である．より正確には，
小拡大$R'\twoheadrightarrow R$の核を$I$とし，$R$上の根付き変形とその同定を固定すると，
その持上げの同型類集合は非空で，
\[
 \HH^2(X)^{s_{\HH}}\otimes_{\C}I
\]
のtorsorとなる．従って任意の不変な一次方向は，冪関係を保った互換的な全Artin次数の変形へ延長できる．
\end{theorem}
\begin{proof}
通常のAbel変形理論の障害は$\HH^3(X)\otimes I$に属するので，
圏自体の持上げは存在する．$R$上の同定を固定した持上げ集合$\mathcal L$は
$W=\HH^2(X)\otimes I$のaffine torsorである\cite{LVdB,BW}．
ここでは代数の厳密な基底付き積ではなく，特殊化の同定を保つMorita/Abel同値を用いる．
Tiltingによる両記述の互換性は\cite[Theorem 3.18, Proposition 3.21]{BW}，
導来核のHochschild作用との互換性は\cite[Theorem 3.2, \S4.6]{KellerDIH}による．

$s$による特殊ファイバーの再標識と，指定した$s_R$による$R$上の同定を組み合わせると，
$\mathcal L$上のaffine自己写像$a$を得る．その線型部分は$s_{\HH}\otimes1_I$である．
$a$の固定点は，その変形上へ$s_R$が関手として持ち上がることに等しい．
これは変形groupoidに自己同値をtransportする定義そのものであり，
接空間の固有値から高次の持上げを推測しているのではない．

重要なのは$a^r=1$が\emph{実際の持上げtorsor上で}成立することである．
$s_R^r\simeq\tau_R$であるため，$a^r$は$\tau_R$による再標識となる．
任意の持上げ$\mathcal A_{R'}$は命題\ref{prop:tilting-serre}の$\tau_{R'}$を持ち，
これは$\tau_R$を持ち上げる．従ってこの再標識は$\mathcal L$の同型類を変えない．
ここで相対Serreの存在を使っており，$s_{\HH}^r=1$だけから$a^r=1$を結論してはいない．

任意の$x\in\mathcal L$に対して，affine空間での重心
\[
 \bar x=\frac1r\sum_{j=0}^{r-1}a^j(x)
\]
は$a$の固定点である．係数の和が1なので原点の選択に依存しない．
固定点全体は$W^{s_{\HH}}$のtorsorになる．固定点で得た関手の持上げは
定理\ref{thm:artin-power}により冪関係も満たし，補題\ref{lem:artin-heart}によりexactである．
$\HH^1=0$による核の一意性から，根の持上げを加えても新しいtorsor方向は生じない．
これを小拡大の列に沿って反復する．

冪関係の自然同型を固定しても，その特殊化への持上げの曖昧さはscalarである．
$1+\mathfrak m_R$上の$r$乗写像はlogarithmにより同型なので，
scalar自己同型で補正でき，新しい無限小パラメータは増えない．
この主張は根の同型類についてであり，scalarの慣性を消した2-stack同値を主張していない．
\end{proof}

\begin{remark}[形式性・巡回性・代数化の区別]\label{rem:formal-boundaries}
定理\ref{thm:formal-roots}は任意の小拡大についての持上げ定理である．
完成局所基底上では互換的なArtin切断の系を得るが，有限型代数化や一般ファイバーの指数の分類は含まない．
$s$自身の有限位数作用，cochainレベルの厳密な$C_r$作用，Haniharaの巡回不変root pairも構成していない．
有限性を証明したのは，相対Serreを自動的に持ち上げられる\emph{変形の持上げ集合}への作用である．
また$\HH^1\ne0$では核の一意性を使えず，本定理をそのまま適用しない．
\end{remark}

\subsection{反標準豊富性と例外列}
\begin{proposition}[Artin変形での反標準ample pair]\label{prop:relative-ample}
$X$を命題\ref{prop:tilting-serre}のFano多様体とし，$\mathcal O_R$を構造層のflatな持上げとする．
この持上げは存在し同型まで一意で，$(\mathcal O_R,\tau_R)$は$R$線形Artin--Zhang ample pairである．
任意の$F\in\mathcal A_R$について
\[
 \Ext^{q}_{\mathcal A_R}(\tau_R^{-m}\mathcal O_R,F)=0
 \quad(q>0,\ m\gg0)
\]
が成立する．$X$が次数$d$のdel Pezzo曲面なら，任意の$m\ge0$について
$\operatorname{Hom}(\mathcal O_R,\tau_R^m\mathcal O_R)$は階数
$1+d m(m+1)/2$の自由$R$加群である．
\end{proposition}
\begin{proof}
$H^{>0}(\mathcal O_X)=0$なので，Lowenの対象持上げ理論から$\mathcal O_R$の存在と一意性を得る．
$E_m=\tau_R^{-m}\mathcal O_R$は$R$-flatであり，特殊ファイバーのpushforward $i_*G$について
導来随伴により
\[
 \mathbf R\operatorname{Hom}_{\mathcal A_R}(E_m,i_*G)
 \simeq\mathbf R\operatorname{Hom}_X(\omega_X^m,G)
\]
となる．$F$の有限$\mathfrak m_R$-filtrationの各subquotientにSerre消滅を適用すると，
表示の漸近消滅を得る．全射の核に$q=1$の消滅を適用すればAZの漸近全射条件となる．
特殊ファイバーの$F/\mathfrak m_RF$を，十分負に捩った有限個の構造層で生成する．
その射は$\mathfrak m_RF$への$\Ext^1$消滅により$F$へ持ち上がる．
余核は$\mathfrak m_R$を法として零なのでnilpotent Nakayamaにより零であり，生成条件も従う．

最後に導来Hom複体
\[
 \mathbf R\operatorname{Hom}(\mathcal O_R,\tau_R^m\mathcal O_R)
\]
は$R$-perfectである．
特殊ファイバーはKodaira消滅とRiemann--Rochから次数0だけにあり，次元は表示の値である．
局所Artin環上のperfect複体の最小自由モデルを取ると，特殊化の消滅から次数0以外の項を消去できる．
従って自由加群としての主張を得る．
\end{proof}

\begin{proposition}[Full例外列の相対持上げ]\label{prop:relative-exceptionals}
上のtilting設定で，$\Db(\coh X)$のfull exceptional collection $(E_1,\ldots,E_l)$は
任意のArtin変形の相対perfect圏へ持ち上がる．各対象の持上げは同型まで一意で，
その自己導来準同型は$R$，後向きの導来Homは零であり，列は相対perfect圏を生成する．
根が持ち上がっており，特殊ファイバーの列の項が$s$の軌道に属するなら，その軌道関係も持ち上がる．
\end{proposition}
\begin{proof}
各対象の$\Ext^1,\Ext^2$の消滅から存在と一意性を得る．
相対Homは$R$-perfectであり，特殊化は自己Homの場合$\C$，後向きには零なので，
derived Nakayamaと最小自由モデルから所望の相対Homを得る．
相対evaluationを用いて有限個の左mutationを行えば，任意の相対perfect対象から右直交部分を取り出せる．
その特殊化は元の列のfull性で零となり，derived Nakayamaで対象自身が零となる．
軌道関係は，同一の例外対象の二つの持上げの一意性から従う．
\end{proof}

これは$R$-linearな例外列と$\Perf(R)$による相対生成である．
$R$が非reducedの場合の全有界導来圏を$\C$上exceptionalな有限列で生成すると主張したものではない．
生成付き根と残余mutationへの道筋を与えるが，root pairの巡回不変性はなお別条件である．

\section{Fano多様体のtensor根と曲面のPoisson方向}
\label{sec:surface}
$X$を滑らかな複素Fano多様体とする．HKR分解は
\begin{equation}\label{eq:hkr}
 \HH^2(X)=U\oplus P,\qquad U=H^1(X,T_X),\quad P=H^0(X,\wedge^2T_X)
\end{equation}
となる．$H^2(\OO_X)=0$を使った．曲面では$P=H^0(-K_X)$で，全双ベクトルはPoissonである．
高次元の双ベクトルは一次変形を定めるが，二次以上にはJacobi障害があり得る．$P$の作用は$\wedge^2T_X$の\emph{自然な線型化}で取り，
反標準写像の射影作用だけで置き換えない．

\begin{lemma}[三角作用]\label{lem:triangular}
$s=L\otimes f^*$の$\HH^2$作用は，\eqref{eq:hkr}について
\[
 \alpha(u,\pi)=(f^*u+B_L\pi,\ f^*\pi)
\]
という形である．ここで$B_L:P\to U$はある線型写像である．
\end{lemma}
\begin{proof}
線束のtensorは局所的に恒等Morita同値なので，局所的な非可換変形の記号である
大域Poisson成分を変えない．一方，$H^2(\OO_X)=0$により$L$は任意の一次可換変形へ持ち上がり，
そのtensor自己同値は同じ可換変形上にある．従ってtensorは$U$を恒等に保つ．
$f^*$は通常のHKR分解と互換である．以上から表示を得る．
\end{proof}

\begin{theorem}[通常の指数のtensor根は一次変形を制限しない]\label{thm:tensor-root}
滑らかな複素Fano多様体で$L^{\otimes r}\simeq\omega_X^{-1}$なら，
$T_L=(-)\otimes L$は$\HH^2(X)$に恒等に作用する．従って任意の一次Abel変形に
$T_L$はexact自己同値として持ち上がる．
さらに$s=L\otimes f^*$が根であるとき，その一次持上げ条件は単に
$f_{\HH}(u)=u$である．相対Serre仮定の下では，これらは根の冪関係も満たす持上げになる．
\end{theorem}
\begin{proof}
補題\ref{lem:triangular}で$f=1$とすれば，$N=(T_L)_{\HH}-1$は$P$から$U$へ写し，$U$を消すので$N^2=0$．
$(T_L)^r=\tau$と命題\ref{prop:finite-hh}より$(1+N)^r=1$だから$rN=0$，従って$N=0$である．
定理\ref{thm:fixed}と系\ref{cor:relative-root}を適用する．最後の条件では
$s_{\HH}=(T_L)_{\HH}f_{\HH}=f_{\HH}$である．
\end{proof}

\begin{proposition}[不変Poisson方向の補償]\label{prop:invariant-poisson}
$s^r=\tau$なら，以下は完全列である：
\begin{equation}\label{eq:invariants}
 0\longrightarrow U^{f^*}\longrightarrow\HH^2(X)^{s_{\HH}}
 \longrightarrow P^{f^*}\longrightarrow0.
\end{equation}
さらに$f^*|_U=1$なら，HKRの\emph{元の分解}において
\begin{equation}\label{eq:trivial-U}
 \HH^2(X)^{s_{\HH}}=U\oplus P^{f^*}.
\end{equation}
\end{proposition}
\begin{proof}
$0\to U\to\HH^2\to P\to0$は命題\ref{prop:finite-hh}により有限巡回群の表現の完全列である．
標数零の不変部分関手は平均化によりexactなので\eqref{eq:invariants}を得る．
$f^*|_U=1$とし$\pi\in P^{f^*}$を取ると，補題\ref{lem:triangular}より
$\alpha^r(0,\pi)=(rB_L\pi,\pi)$．$\alpha^r=1$だから$B_L\pi=0$である．
従って$(0,\pi)$自身が不変となり，\eqref{eq:trivial-U}が従う．
\end{proof}

\begin{remark}[Picard re-markingとの接続]\label{rem:picard}
局所積を$ab+\varepsilon\pi(da,db)/2$と規格化し，線束の遷移関数を$g_{ij}$とすると，
そのstar共役は$\id+\varepsilon\pi^\sharp(d\log g_{ij})$である．
従ってtensorによるPoisson成分から可換成分への混合は，向きの規約を固定すれば
\[
 \delta_\pi(L)=H^1(\pi^\sharp)c_1(L)\in H^1(T_X)
\]
で与えられる．符号を反転してもその核は変わらない．
これは別研究\cite{Marking}のmarking忘却の障害と同じ一次機構であり，
\emph{線束を対象として持ち上げること}と\emph{tensorを固定変形の自己同値として持ち上げること}は異なる．
Bursztyn--Dolgushev--Waldmann\cite[Theorem 3.11]{BDW}は滑らかな実多様体上のstar productで
Picard--Morita作用を$B$-field作用と同定する．その$C^\infty$形式理論を，本稿の射影代数的変形へ無条件に転用しない．
\end{remark}

\begin{proposition}[圏自身の形式的非障害性]\label{prop:hh3}
滑らかな複素del Pezzo曲面では$\HH^3(X)=0$である．
従ってLowen--Van den BerghのAbel変形関手では，小拡大に関する障害が消える．
これ単独では根の持上げを意味しないが，$\HH^1=0$のとき定理\ref{thm:formal-roots}と組み合わせられる．
相対Fano性は命題\ref{prop:relative-ample}の別の議論を使い，代数化はなお含まない．
\end{proposition}
\begin{proof}
HKRにより$\HH^3=H^1(-K_X)\oplus H^2(T_X)$．
$-K_X=K_X+(-2K_X)$とKodaira消滅から第一項は零．
Serre双対性から第二項の双対は$H^0(\Omega_X^1\otimes K_X)$である．
非零反標準切断を掛ける単射$\Omega_X^1\otimes K_X\hookrightarrow\Omega_X^1$と
合理曲面の$H^0(\Omega_X^1)=0$によりこれも零．
変形の障害は\cite[Theorem 3.1]{LVdB}を用いる．
\end{proof}


\section{次数2：Geiser根を消す全Poisson方向}
\label{sec:degree2}
$X$を次数2の滑らかなdel Pezzo曲面とする．反標準二重被覆の表示は
\[
 X=\{w^2=F_4(x_0,x_1,x_2)\}\subset\mathbb P(1,1,1,2)
\]
で，Geiser involutionは$\gamma(w)=-w$である\cite[\S2.3, \S4]{DM}．
任意の$(-1)$曲線$E$について$E+\gamma E\sim-K_X$なので，
$s_E=\OO_X(E)\otimes\gamma^*$は二乗根である\cite{Index}．

\begin{lemma}[自然な反標準表現]\label{lem:geiser-character}
$\gamma$は$U=H^1(T_X)$に恒等，$P=H^0(\wedge^2T_X)$に$-1$として作用する．
$\dim U=6$，$\dim P=3$である．
\end{lemma}
\begin{proof}
反標準二重被覆の一次変形はbranch quarticの変形で与えられる．
実際，Kodaira消滅で反標準環の各正次数部分はbase changeし，
次数1の三生成元と次数2の生成元の関係式は，標数零で平方完成して同じ表示にできる．
$\gamma$は全branch変形上に作用し基底を固定するので，$U$上恒等となる．
四次式の変形は射影次元14，$\operatorname{PGL}_3$は次元8で，滑らかな四次曲線の自己同型は有限だから$\dim U=6$．

局所座標$u,v$においてcanonicalの生成元を$du\wedge dv/w$と書ける．
$\gamma$はこれを負にする．その逆のbivectorも負にし，反標準切断の係数となる
$x_0,x_1,x_2$は固定する．従って自然な反標準表現は$-I_3$である．
分岐点でもresidueによる同定は正則なので，大域的にこの指標となる．
\end{proof}

\begin{theorem}[Geiser根の一次持上げ障害]\label{thm:geiser}
任意の$L$で$L\otimes\gamma^*L\simeq\omega_X^{-1}$となるものについて，
$s=L\otimes\gamma^*$が自己同値として持ち上がる一次変形空間は
\[
 \HH^2(X)^{s_{\HH}}=H^1(T_X)\oplus0
\]
である．特に非零Poisson成分を持つ変形には$s$は持ち上がらない．
さらに$\Aut(X)=\{1,\gamma\}$なら，非零Poisson成分を持つ一次変形上には，
特殊化で非自明な根となるexact相対反Serre根は存在しない．
\end{theorem}
\begin{proof}
最初の主張は命題\ref{prop:invariant-poisson}と補題\ref{lem:geiser-character}による．
後半を確認する．任意の古典的exact根$t=M\otimes f^*$で$t^r=\tau$なら，
norm条件と$K_X$との交点から$r\mid K_X^2=2$である．
$r>1$なら$r=2$．$f=1$では$2c_1(M)=-K_X$となるが，$(-1)$曲線と交わらせると
偶数$=1$となり不可能．従って$f=\gamma$であり，前半の障害を適用できる．
相対根は特殊化すると同じ次数の古典的根になるので矛盾する．
\end{proof}

\begin{remark}
滑らかなplane quarticが自己同型を持たない一般の場合には$\Aut(X)=\{1,\gamma\}$である\cite{DM}．
$\gamma$は全ての反標準divisorを集合として保つが，そのPoisson切断は負にする．
従って零divisorの不変性だけでは固定変形への持上げを保証できない．
定理\ref{thm:geiser}は一次障害であり，次の系は全形式次数へ強める．ただし特定の代数化された非可換曲面の
全自己同値を分類したという主張ではない．
\end{remark}

\begin{corollary}[Geiser根を保つ全形式変形は可換]\label{cor:geiser-formal}
次数2の任意の滑らかな複素del Pezzo曲面$X$と，Geiser型二乗根$s=L\otimes\gamma^*$を固定する．
標識付きflat Artin Abel変形で$s$を相対二乗根として保つものの同型類は，
$X$の可換な幾何学的Artin変形の同型類と自然に一対一に対応する．
この根付き変形問題は形式的に滑らかで一次次元6である．
従って一次項でPoisson成分が零という条件だけでなく，\emph{全形式次数で幾何学的変形に同値}となる．
\end{corollary}
\begin{proof}
合理曲面はHille--Perling\cite[Theorem 1.1]{HP}によりtilting bundleを持つ．
次数2の自己同型群は有限なので$H^0(T_X)=0$，また$H^1(\OO_X)=0$であり$\HH^1(X)=0$．
命題\ref{prop:hh3}と定理\ref{thm:formal-roots}から，根付き変形は非障害的で接空間は$H^1(T_X)$である．

幾何学的変形も$H^2(T_X)=0$により非障害的である．
反標準二重被覆とそのdeck involutionはArtin変形上で定義されるので$\gamma$が持ち上がる．
$H^1(\OO_X)=H^2(\OO_X)=0$により，$L$は幾何学的変形へ一意に持ち上がる．
同じ消滅からPicard群の特殊化の核は零なので，
$L_R\otimes\gamma_R^*L_R=\omega_{X_R/R}^{-1}$も成立する．従って幾何学的変形から根付きAbel変形への写像がある．

この写像が全Artin基底で全単射であることを長さに関する帰納法で示す．
小拡大の両側の持上げはそれぞれ$H^1(T_X)\otimes I$と$\HH^2(X)^{s_{\HH}}\otimes I$のtorsorである．
定理\ref{thm:geiser}によりこれらは同じ空間で，写像の線型部分もHKRの標準包含すなわち同型である．
従って持上げtorsor間の写像は全単射となる．特殊ファイバーから帰納して結論を得る．
根の自然同型のscalar慣性を保持したstack同型は主張していない．
\end{proof}

\section{次数4：十個の座標部分空間}
\label{sec:degree4}
$X$を滑らかな次数4 del Pezzo曲面とし，
\begin{equation}\label{eq:diagonal}
 X=\left\{\sum_{i=0}^4x_i^2=\sum_{i=0}^4\lambda_i x_i^2=0\right\}\subset\mathbb P^4,
 \qquad \lambda_i\ne\lambda_j\ (i\ne j)
\end{equation}
と書く．符号変換群$G=(\{\pm1\}^5)/\{\pm I\}\simeq(\Z/2)^4$が作用する．
一般の$X$では$\Aut(X)=G$である\cite[Introduction, \S2]{HKT}．

\begin{lemma}[符号群の表現と二乗根]\label{lem:degree4-action}
$U=H^1(T_X)$は2次元の自明$G$表現である．
$P=H^0(\wedge^2T_X)$を座標$x_0,\ldots,x_4$で表示すると，
符号変換$\epsilon=(\epsilon_0,\ldots,\epsilon_4)$の自然な作用は
\begin{equation}\label{eq:det-character}
 x_i\longmapsto\left(\prod_{j=0}^4\epsilon_j\right)^{-1}\epsilon_i x_i.
\end{equation}
各$1\ne f\in G$について，$L\otimes f^*L\simeq\omega_X^{-1}$となる線束$L$が存在する．
\end{lemma}
\begin{proof}
Euler列と$H^1(\OO_X(1))=H^2(\OO_X)=0$から$H^1(T_{\mathbb P^4}|_X)=0$．
従って全一次変形は二つの二次式の変形から得られる．
固有値が相異なるので，dual numbers上でもpencilを同時対角化できる．
五つの$\lambda_i$の変形を$\operatorname{PGL}_2$で割った2方向を符号群は固定する．
また$h^0(N_{X/\mathbb P^4})=2(15-2)=26$，$h^0(T_{\mathbb P^4}|_X)=24$，
$h^0(T_X)=0$から$\dim U=2$を得る．

自然なcanonical線型化をresidueで計算する．$\Omega=\sum(-1)^i x_i\,dx_0\wedge\cdots\widehat{dx_i}\cdots\wedge dx_4$
は$\det(\epsilon)$倍され，二つの方程式は固定される．
従ってadjunctionのcanonical指標は$\det(\epsilon)$で，反標準切断にはその逆が掛かる．
これが\eqref{eq:det-character}である．$\epsilon$を$-\epsilon$に変えても作用は同じなので$G$上well-definedである．

非自明な符号変換は負符号の集合$S$が大きさ1または2の代表を唯一持つ．
$|S|=1$なら$j\notin S$，$|S|=2$なら$j\in S$を取る．
pencilのrank4 quadric $Q_j=\{\sum_{i\ne j}(\lambda_i-\lambda_j)x_i^2=0\}$の
四次元非退化部分で$f$のdeterminantは$-1$である．従ってsmooth quadric surfaceの
二つのrulingを交換する．$X$への二つのconic pencilの類を$A_j,B_j$とすれば
$A_j+B_j=-K_X$かつ$f^*A_j=B_j$．$L=\OO_X(A_j)$で主張を得る．
\end{proof}

\begin{theorem}[一次持上げ部分空間の配置]\label{thm:degree4-arrangement}
\eqref{eq:diagonal}の$X$が$\Aut(X)=G$を満たすとする．
非自明な反Serre根のうち少なくとも一つを\emph{関手として}持ち上げられる一次変形の集合は
\begin{equation}\label{eq:arrangement}
 \mathcal L_X=U\oplus\bigcup_{0\le i<j\le4}\{(a_0,\ldots,a_4)\in\C^5:a_i=a_j=0\}
 \subset\HH^2(X)=\C^2\oplus\C^5.
\end{equation}
これは10個の5次元線型部分空間の和集合であり，各成分の余次元は2である．
Poisson部分の被約配置のイデアルは，五つの平方因子を持たない次数4単項式で生成され，Hilbert級数は
\[
 \frac{1+2t+3t^2+4t^3}{(1-t)^3}
\]
である．
この$X$では$\HH^1=0$かつtilting bundleが存在するため，定理\ref{thm:artin-power}から
$\mathcal L_X$は冪関係も満たす根の一次持上げ集合である．
特にPoisson成分の非零座標が4個以上なら相対根の存在を妨げる．
\end{theorem}
\begin{proof}
古典的根$t=M\otimes f^*$の次数$r$は$4$を割る．$G$のexponentは2なので，
$r=4$ならnormは$2(M+f^*M)=-K_X$となり，$(-1)$曲線との交点から不可能．
$r=2,f=1$も同じ理由で不可能．従って非自明な根は$1\ne f\in G$による二乗根だけであり，
各$f$は前補題で実現する．

$|S|=1$なら\eqref{eq:det-character}の$+1$固有空間はその一座標軸．
$|S|=2$なら$+1$固有空間は$S$に属する二座標を零とする3次元空間である．
$U$上の作用は恒等なので，命題\ref{prop:invariant-poisson}により，
\emph{同じ$f$を持つ全ての線束解$M$について}固定部分は$U\oplus P^{f^*}$となる．
五つの座標軸は十個の3次元空間の和集合に既に含まれるから\eqref{eq:arrangement}を得る．

配置に消える単項式は少なくとも4変数を含むものに限る．よって被約イデアルは
$\bigcap_{i<j}(a_i,a_j)=(\prod_{k\ne i}a_k:0\le i\le4)$である．
残る単項式のsupportを数えるとHilbert級数は
\[
 1+5\frac{t}{1-t}+10\left(\frac{t}{1-t}\right)^2+
 10\left(\frac{t}{1-t}\right)^3,
\]
すなわち表示の有理式となる．最後の主張には定理\ref{thm:artin-power}を適用する．
\end{proof}

\begin{remark}[scheme構造を主張しない]
$\mathcal L_X$は様々な標識根に対する一次持上げ空間の\emph{集合論的和集合}である．
その被約座標配置を計算したのであって，高次の根モジュライの接錐・scheme構造を同定したのではない．
$\HH^3=0$による圏自身の形式的非障害性も，この和集合の全Artin次数での線型性を意味しない．
\end{remark}

\begin{corollary}[各符号根の全次数持上げ]\label{cor:degree4-formal}
次数4の曲面と非自明な符号変換$f$を取り，$s=L\otimes f^*$を二乗根とする．
$s$を保つArtin変形問題は形式的に滑らかで，$|S|=1$なら一次次元3，$|S|=2$なら一次次元5である．
それぞれの不変一次方向は全形式次数へ，相対ample pairと根の冪関係を保って延長できる．
\end{corollary}
\begin{proof}
$\HH^1(X)=0$，命題\ref{prop:hh3}，合理曲面のtilting bundleを用いて定理\ref{thm:formal-roots}を適用する．
次元は$2+\dim P^{f^*}$であり，補題\ref{lem:degree4-action}から$2+1$または$2+3$となる．
相対ample性は命題\ref{prop:relative-ample}による．
\end{proof}

\begin{example}[同じ接空間から同じ固定形式部分は従わない]\label{ex:nonlinear-fixed}
$\C[[x,y]]$の二つのinvolution
\[
 a(x,y)=(x,-y),\qquad b(x,y)=(x,2x^2-y)
\]
は同じ一次作用を持つが，固定形式部分はそれぞれ$y=0$と$y=x^2$である．
合成$b\circ a$は$y$を$y+2x^2$へ移し無限位数となる．
従って各根の固定変形問題が滑らかであっても，異なる線束解・異なる根を一緒に扱うと，
接空間の配置だけで高次の和集合や同時線型化は決まらない．
この例はdel Pezzoの根の実現ではなく，そのような推論の形式幾何上の障害である．
\end{example}

\section{Fermat三次曲面：二次元の根付き形式変形}
\label{sec:fermat}
\begin{proposition}\label{prop:fermat-formal}
$X=\{x^3+y^3+z^3+w^3=0\}\subset\mathbb P^3$，$\zeta^3=1$，$\zeta\ne1$とし，
$f[x:y:z:w]=[x:y:z:\zeta w]$，$L=\OO_X(\{x+y=z+w=0\})$と置く．
$s=L\otimes f^*$は三乗根であり，$\HH^2(X)$上の固有値$1,\zeta,\zeta^2$の重複度は$2,3,3$である．
従って$s$を保つArtin変形は形式的に滑らかで一次次元2となる．
その不変空間は，可換な1方向と非零Poisson成分へ写る1方向を持つ．
\end{proposition}
\begin{proof}
三つの巡回直線の和は平面$x+y=0$の切断なので，normは$-K_X$となる．
従って$s^3=\tau$である．$H^1(T_X)$はJacobian環の次数3部分であり，
基底を$xyz,xyw,xzw,yzw$と取れる．$f$の重みは$0,1,1,1\pmod3$である．
一方$H^0(-K_X)$の自然表現には線型代表のdeterminantの逆が掛かるので，
$x,y,z,w$の重みは$2,2,2,0$となる．
線束tensorの混合は三角部分だけで，有限位数作用はsemisimpleである．
よって合計した重複度は$2,3,3$となり，命題\ref{prop:invariant-poisson}から
不変空間に$H^1(T_X)^{f^*}$と$H^0(-K_X)^{f^*}=\C w$がそれぞれ一方向ずつ現れる．

三次曲面の自己同型群は有限なので$\HH^1=0$であり，$\HH^3=0$とtilting bundleの存在を合わせて
定理\ref{thm:formal-roots}を適用する．$w$へ写る不変類を取れば非可換な一次方向を持つ形式族が得られる．
HKR分解で$(0,w)$自体が不変とは主張しない．$H^1(T_X)$方向の補償が必要な場合でも，
平均化で不変なliftを取れる．
\end{proof}

\section{同じ数値根を区別する射影平面の例}
\label{sec:plane}
\begin{proposition}\label{prop:plane}
$X=\mathbb P^2$，$\zeta^3=1,\zeta\ne1$，
$f[x:y:z]=[x:\zeta y:\zeta^2z]$とし，
$s_0=(-)\otimes\OO(1)$，$s_1=\OO(1)\otimes f^*$とする．
両者は同じ数値作用を持つ正規化三乗根であるが，
\[
 \dim\HH^2(X)^{s_0}=10,\qquad \dim\HH^2(X)^{s_1}=4.
\]
後者の不変Poisson切断は$\langle x^3,y^3,z^3,xyz\rangle$である．
\end{proposition}
\begin{proof}
$\HH^2(X)=H^0(\OO(3))$で$H^1(T_X)=0$だから線束tensorの作用は恒等である．
$f$は数値群上恒等なので$s_0,s_1$の数値作用は等しく，$f^3=1$より両者とも三乗根である．
$f$の線型代表のdeterminantは1であり，反標準自然作用は三次式への代入に一致する．
$x^ay^bz^c$の重みは$b+2c\pmod3$．$a+b+c=3$で重み零の単項式は表示の四つだけである．
他の二固有空間の次元はそれぞれ3となる．
\end{proof}

この例では指定した$s_1$が失われる変形でも$s_0$は\emph{関手として}持ち上がる．
従って一つの根の障害を，その圏の全ての根の非存在と同一視しない．
次数2・4で全根を扱えたのは，古典的自己同値群とnorm解を別に制御したためである．


\section{有限な反Serre商と複数の根の同時持上げ}
\label{sec:simultaneous}
一つの根の持上げと，複数の指定された根を同一変形上で保つことは異なる．
本節では自己同値の\emph{同型類の群}を用いる．群作用のcoherentな自然同型は標識しない．

\begin{theorem}[有限反Serre商による同時持上げ]\label{thm:finite-package}
定理\ref{thm:formal-roots}の仮定，すなわちtilting bundleと
$\HH^1(X)=\HH^3(X)=0$を持つ滑らかな連結射影$X$を考える．
$\Gamma\subset\Aut_{\mathrm{ex}}(\coh X)$を$\tau$を含む有限生成部分群とし，
\[
 H=\Gamma/\langle\tau\rangle
\]
が有限であると仮定する．$\Gamma$の指定した生成元を同時に持ち上げる
flat Artin Abel変形問題は形式的に滑らかで，一次空間は$\HH^2(X)^\Gamma$である．
小拡大の核$I$に対する持上げの同型類は
$\HH^2(X)^\Gamma\otimes I$の非空torsorとなる．
特殊ファイバーの群関係は，持ち上げた関手の\emph{同型類の間の関係}として保たれる．
特に生成元が反Serre根なら，その冪関係も保たれる．
\end{theorem}
\begin{proof}
既存の根付き変形を$R$上で固定する．$\HH^3=0$により，$R'$上への圏自体の持上げ集合
$\mathcal L$は$W=\HH^2(X)\otimes I$の非空affine torsorである．
$\HH^1=0$により，固定変形上の各可逆核の持上げは同型まで一意である．
従って生成元の合成の関係は，特殊化で同じ可逆核を与える関手の同型として保たれる．
このため，再標識は$\mathcal L$上に$\Gamma$の整合した作用を定める．
命題\ref{prop:tilting-serre}の相対$\tau_{R'}$が全持上げに存在するので，
$\tau$の再標識は恒等となり，作用は有限群$H$を経由する．
$W$上の線型部分も$H$を経由する（命題\ref{prop:finite-hh}のSerre自明性）．

任意の$x\in\mathcal L$の軌道のaffine重心
$|H|^{-1}\sum_{h\in H}h(x)$は共通固定点である．
固定点の差はちょうど$W^H$であり，共通固定点は全生成元の持上げを与える．
heart-exact性と群関係はそれぞれ補題\ref{lem:artin-heart}，
補題\ref{lem:kernel-rigid}から従う．小拡大を反復して結論を得る．
群関係を実現する自然同型の全coherenceを選択したという主張ではない．
\end{proof}

\begin{remark}[有限性の位置]
必要なのは$\Gamma/\langle\tau\rangle$の\emph{関手同型類としての有限性}である．
各生成元のHH作用が有限位数という条件だけでは十分でない．
以下ではこの有限性を，Picard格子の整性と可換性から実際に判定する．
\end{remark}

\section{次数4の同時二乗根と二部グラフ}
\label{sec:root-graphs}
次数4の滑らかなdel Pezzo曲面を節\ref{sec:degree4}の対角表示で考える．
追加の自己同型が存在する特殊な曲面でも，以下の指定した符号変換部分群に関する主張は成立する．
五つの点をblowupした表示により
\[
 \Pic X=\Z h\oplus\bigoplus_{i=0}^4\Z e_i,\quad
 K=-3h+\sum_i e_i,\quad A_i=h-e_i,\quad v_i=A_i+K/2
\]
と置く．五つのconic pencilの対を，pencilの五つの特異二次曲面とこの添字で対応させる．
既知の$D_5$格子・conic pencilの幾何\cite[\S2]{HKT}を用いるが，以下の整性と群計算は直接行う．
\[
 K^2=4,\qquad K\cdot v_i=0,\qquad v_i\cdot v_j=-\delta_{ij}.
\]
$K,v_0,\ldots,v_4$は$\Pic(X)_\Q$の基底であって，整数基底ではない．
$i\ne j$に対し，$f_{ij}$を座標$x_i,x_j$の符号を反転するinvolutionとする．
補題\ref{lem:degree4-action}のruling判定から，$f_{ij}$は$v_i,v_j$を負にし，他の$v_k$と$K$を固定する．

\begin{lemma}[全線束解の整性]\label{lem:graph-norm}
$s=\OO_X(D)\otimes f_{ij}^*$が二乗根となるための必要十分条件は
\begin{equation}\label{eq:edge-root}
 D=-K/2+a_{ij,i}v_i+a_{ij,j}v_j,\qquad
 a_{ij,i},a_{ij,j}\in\Z,\quad a_{ij,i}+a_{ij,j}\equiv1\pmod2
\end{equation}
である．これは有効divisorを仮定しない全ての線束解を記述する．
\end{lemma}
\begin{proof}
$D+f_{ij}^*D=-K$を有理基底で解くと，表示の形となる．
その$h,e_0,\ldots,e_4$成分は
\[
 \begin{aligned}
 [h]D&=\frac{3-a-b}{2},&
 [e_i]D&=\frac{-1-a+b}{2},\\
 [e_j]D&=\frac{-1+a-b}{2},&
 [e_k]D&=\frac{-1+a+b}{2}\quad(k\notin\{i,j\}).
 \end{aligned}
\]
$e_k$成分と$e_i,e_j$成分の差から$a,b\in\Z$が必要となり，
残る条件は$a+b$が奇数ということである．逆は同じ式から直ちに従う．
\end{proof}

$\mathcal G=(V,E)$を$V=\{0,1,2,3,4\}$上の非空な単純グラフとする．
辺$e=\{i,j\}$ごとに\eqref{eq:edge-root}の根$s_e$を選ぶ．
$\Gamma_E=\langle s_e:e\in E\rangle$と置くと，各$s_e^2=\tau$だから$\tau\in\Gamma_E$である．
$V_E$は辺に接する頂点集合，$c_E$は$(V_E,E)$の連結成分数とする．

\begin{theorem}[二部性による可換な根の選択]\label{thm:bipartite-roots}
次の三条件は同値である．
\begin{enumerate}
\item $\mathcal G$は二部グラフである．
\item 各辺の線束を選び，$s_es_d\simeq s_ds_e$が全ての辺対について成立するようにできる．
\item 各辺の線束を選び，$\Gamma_E/\langle\tau\rangle$が有限となるようにできる．
\end{enumerate}
さらに\emph{既に選んだ}根が可換となるための必要十分条件は，同じ頂点に接する全ての辺で
その頂点係数$a_{e,i}$が一致することである．この場合$a_i=a_{e,i}$と書け，
\begin{equation}\label{eq:color-parity}
 a_i\in\Z,\qquad a_i+a_j\equiv1\pmod2\quad(\{i,j\}\in E)
\end{equation}
となる．この整係数の選択が可換な根の全てを記述し，
\[
 \Gamma_E/\langle\tau\rangle\simeq(\Z/2)^{|V_E|-c_E}
\]
である．
\end{theorem}
\begin{proof}
線束tensorとpullbackの半直積表示を用いる．$f_e,f_d$は可換なので
\begin{equation}\label{eq:root-commutator}
 [s_e,s_d]=(-)\otimes\OO_X(C_{ed}),\qquad
 C_{ed}=(1-f_d^*)D_e-(1-f_e^*)D_d.
\end{equation}
辺が交わらなければ$C_{ed}=0$．共有頂点が$i$なら
\[
 C_{ed}=2(a_{e,i}-a_{d,i})v_i.
\]
Picard群はtorsion-freeであり，$v_i\ne0$なので，可換性は係数の一致と同値となる．
補題\ref{lem:graph-norm}により，その共通係数は\eqref{eq:color-parity}を満たす．
$a_i\bmod2$が二彩色を与えるから(ii)は(i)を含意する．
逆に二彩色$a_i\in\{0,1\}$を選べば，辺の線束は色1の端点のconic pencil
$\OO_X(A_i)$となり，所望の可換な根を得る．

(ii)なら商は可換な位数2の有限個の元で生成されるので(iii)となる．
逆に商が有限なら，\eqref{eq:root-commutator}のある正の冪は$\tau$の冪である．
したがって$nC_{ed}=-mK$となる．$K\cdot C_{ed}=0$と$K^2=4$から$m=0$，
Picard群のtorsion-free性から$C_{ed}=0$となり(ii)を得る．

最後に共通係数を用いて$B=\tfrac12\sum_{i\in V_E}a_i v_i\in\Pic(X)_\Q$と置くと
$D_e=-K/2+(1-f_e^*)B$である．従って$s_e^{\pm1}$からなる任意のwordは，
符号付き長さを$m$，自己同型部分を$F$として
\[
 \OO_X\bigl(-mK/2+(1-F^*)B\bigr)\otimes F^*
\]
と書ける．$F=1$なら$-mK/2$が整数Picard類であり，$-K$のprimitive性から$m$は偶数である．
このwordは$\tau^{m/2}$に等しい．よって商は$f_e$が生成する符号群に単射で，かつ全射である．
その$\mathbb F_2$上のrankはグラフのincidence行列のrank $|V_E|-c_E$である．
各連結成分のspanning treeの辺が偶数和部分空間の基底を与えることから，このrank公式が従う．
$B$は計算用の\emph{有理}Picard類であり，実際の線束による共役を構成したとは主張していない．
\end{proof}

\begin{corollary}[二部グラフの根の同時形式変形]\label{cor:graph-formal}
定理\ref{thm:bipartite-roots}の可換な根を固定する．これらを同時に保つ標識付きflat Artin Abel変形は
形式的に滑らかで，一次空間は，元のHKR分解において
\[
 H^1(T_X)\oplus\operatorname{span}\{x_k:k\notin V_E\}.
\]
その次元は$7-|V_E|$であり，全ての不変一次方向は冪関係と可換性を保って全Artin次数へ延長する．
$V_E=V$なら，この同時根付き変形は全Artin基底で幾何学的な可換変形と同型類で一致する．
\end{corollary}
\begin{proof}
合理曲面のtilting bundle，$\HH^1=\HH^3=0$と有限商を用いて定理\ref{thm:finite-package}を適用する．
各$s_{ij}$の不変部分は線束解によらず
$U\oplus\{a_i=a_j=0\}$である（命題\ref{prop:invariant-poisson}と補題\ref{lem:degree4-action}）．
その共通部分が表示の空間となる．

$V_E=V$のとき接空間は$U=H^1(T_X)$だけである．
全ての幾何学的Artin変形はpencilの固有値の変形で記述され，符号変換が持ち上がる．
$H^1(\OO_X)=H^2(\OO_X)=0$から各線束も一意に持ち上がり，Picard群の関係と根の関係が残る．
したがって幾何学的変形から同時根付きAbel変形への写像がある．
双方の小拡大持上げtorsorが同じ$U\otimes I$で，写像の線型部分は同型なので，
系\ref{cor:geiser-formal}と同じ帰納法で全単射を得る．
\end{proof}

\begin{example}[根の数ではなく被覆頂点数が効く]\label{ex:root-star}
頂点0を中心とする$m$本のstar $E=\{\{0,1\},\ldots,\{0,m\}\}$では，
全ての根の線束を$\OO_X(A_0)$に選べる．$m=1,2,3,4$に対し形式次元はそれぞれ
$5,4,3,2$，Poisson方向の次元は$3,2,1,0$となる．
より少ない三つの根でも，例えば辺$01,12,34$の二部グラフなら全五頂点を覆い，
同時持上げを許す変形は全形式次数で可換に限られる．
これは指定した可換な根の組についての主張で，任意の三根に対する主張ではない．
\end{example}

\begin{remark}[奇閉路が妨げること・妨げないこと]\label{rem:odd-cycle}
三角形などの奇閉路では，どの線束解を選んでも全根を可換にできず，反Serre商も無限となる．
これはPicard格子の奇偶性の障害であり，HH作用だけからは見えない．
しかし各根自体は存在するし，非可換な関手群として全根を同時に変形できないと結論したわけではない．
また二部グラフでも，頂点係数が一致しない線束解を最初から指定した場合は有限商の仮定を満たさない．
従来の10-plane全配置，無限個の線束解，奇閉路についての高次交叉問題は残る．
\end{remark}


\section{奇閉路の具体的な無限商：三次元格子の拡大}\label{sec:triangle-group}
有限反Serre商を失うとき，単に「無限」と言うよりも
Picardの移動格子を記述することが有用である．
頂点$0,1,2$の三角形に対して
\[
 \alpha=\OO_X(A_0)\otimes f_{01}^*,\qquad
 \beta=\OO_X(A_1)\otimes f_{12}^*,\qquad
 \gamma=\OO_X(A_2)\otimes f_{02}^*
\]
を取り，$\Gamma_\triangle=\langle\alpha,\beta,\gamma\rangle$とする．
各根の二乗は$\tau$である．

\begin{proposition}[三角形の反Serre商の平行移動格子]\label{prop:triangle-lattice}
$W=\Q v_0\oplus\Q v_1\oplus\Q v_2$に，
\[
 A=(F_{01},v_0),\quad B=(F_{12},v_1),\quad C=(F_{02},v_2)
 \quad\in\operatorname{Aff}(W)
\]
を作用させる．ここで$F_{ij}$は$v_i,v_j$の符号を反転する線型作用である．
この表示は忠実な群準同型
\[
 \Gamma_\triangle/\langle\tau\rangle\hookrightarrow\operatorname{Aff}(W)
\]
を与え，その像は$A,B,C$の生成群である．
平行移動部分群は
\begin{equation}\label{eq:triangle-translation}
 \Lambda_\triangle
 =\{m_0v_0+m_1v_1+m_2v_2:\ m_i\in\Z,\ 
 m_0\equiv m_1\equiv m_2\pmod2\}
\end{equation}
で，rank $3$である．従って
\[
 1\longrightarrow\Lambda_\triangle
 \longrightarrow\Gamma_\triangle/\langle\tau\rangle
 \longrightarrow(\Z/2)^2\longrightarrow1
\]
が完全である．この完全列の分裂は主張しない．
特に反Serre商は平行移動rank $3$の無限群である．
\end{proposition}
\begin{proof}
任意の$T_Df^*$を，$D$の$W$方向への有理射影と$f^*|_W$からなる
affine変換へ送る．積の公式$T_Df^*T_Eg^*=T_{D+f^*E}(fg)^*$から準同型となり，
$\tau=T_{-K}$は恒等へ送られる．三つの根は$A,B,C$に写る．

逆に，$\Gamma_\triangle$のwordを三つの根の正負冪で書き，符号付き長さを$m$とする．
各根のPicard類は$-K/2+v_i$であり，
全ての符号変換は$K$を固定する．従ってwordの線束部分は
$-mK/2+w$，$w\in W$と書ける．
そのaffine像が恒等なら自己同型部分は恒等で$w=0$．
線束部分が整数Picard類なので，primitiveな$K$により$m$は偶数であり，
元のwordは$\tau^{m/2}$に等しい．よって表示は忠実である．

直接の半直積計算で
\[
 (AB)^2=(I,-2v_1),\quad
 (BC)^2=(I,-2v_2),\quad
 (CA)^2=(I,-2v_0),\qquad
 ABC=(I,v_0-v_1-v_2)
\]
となる．これらが生成する平行移動格子は\eqref{eq:triangle-translation}全体である．
実際，$2v_i$と奇奇奇の平行移動があるので右辺は含まれる．
反対に，$W$の$v_i$座標をmod $2$で見れば，$F_{ij}$は恒等であり，
任意のwordの移動座標のparityは$A,B,C$の出現回数のparityに等しい．
線型部分が恒等なら，三つの出現回数のparityは等しい
（$F_{01}F_{12}F_{02}=1$が唯一の非自明な関係）．
従って全移動座標は同じparityを持ち，逆の包含が従う．
線型部分の像は二つの独立な符号反転で生成されるKlein四元群である．
\end{proof}

\begin{remark}[固定空間と高次の未解決性]
三つの根の共通一次不変空間は，前節のHochschild表現の計算より
\[
 H^1(T_X)\oplus\operatorname{span}\{x_3,x_4\},\qquad\dim=4.
\]
しかし$\Gamma_\triangle/\langle\tau\rangle$は有限ではないため，
定理\ref{thm:finite-package}の平均化は使えない．
この四方向の同時根付き\emph{全Artin持上げ}や二次障害の消滅は未証明である．
反Serre商を一次元平行移動群の拡大として扱う見積りは誤りであり，
実際には\eqref{eq:triangle-translation}の三方向を保持しなければならない．
\end{remark}

\section{明示的な代数族：三つの三乗根を保つ量子平面}
\label{sec:quantum-benchmark}
本節は，形式存在と明示的な代数化を区別するための計算可能な基準例である．
代数族自体とそのAS-regular性・Nakayama公式は既知のskew polynomial理論であり，新しい代数の発見とは扱わない．
また，次数4やFermat三次曲面の形式族をこの例で代数化したことにはしない．

$C=\C[q,q^{-1}]$上で
\begin{equation}\label{eq:quantum-family}
 \mathcal B=C\langle x,y,z\rangle/(yz-qzy,\ zx-qxz,\ xy-qyx),
 \qquad\deg x=\deg y=\deg z=1
\end{equation}
と置く．$\mathcal B$は反復skew polynomial環
\[
 C[x][y;\sigma][z;\rho],\quad
 \sigma(x)=q^{-1}x,\quad\rho(x)=qx,\quad\rho(y)=q^{-1}y
\]
であり，$x^ay^bz^c$が$C$自由基底である．特に全次数部分が有限自由で，これは有限型基底$\mathbb G_m$上の実際の平坦な次数付き代数族である．

\begin{proposition}[複数根を持つ代数的な基準族]\label{prop:quantum-three-roots}
任意の$\lambda\in\C^*$に対し，$A_\lambda=\mathcal B/(q-\lambda)$は
三次元二次Koszul AS-regular代数で，Hilbert級数は$(1-t)^{-3}$，Nakayama自己同型は恒等である．
$\qgr A_\lambda$では$\tau=(3)$となる．
$\zeta^3=1$，$\zeta\ne1$とし，次数付き自己同型
\[
 d(x,y,z)=(x,\zeta y,\zeta^2z),\qquad c(x,y,z)=(y,z,x)
\]
による加群の捩れを$d^*,c^*$と書く．このとき
\[
 s_0=(1),\qquad s_d=(1)d^*,\qquad s_c=(1)c^*
\]
は全てheart-exactで，$s_0^3=s_d^3=s_c^3=\tau$を満たす．
これらは自然同型まで可換で，生成群の$\tau$による商の位数は27を割る．
\end{proposition}
\begin{proof}
PBW基底は上のOre表示から従う．各fiberは三変数skew polynomial環である．
そのquantum Koszul分解の各階数は$1,3,3,1$であり，双対quantum exterior環の頂上次数の積pairingは非退化である．
これによりglobal dimension 3とAS双対性，Gorenstein parameter 3を得る．
同じ事実は通常の多項式環の多重次数付きtwistとしても得られる\cite[Example 5.5]{RRZ}．
同例の規約$x_jx_i=p_{ij}x_ix_j$では
\[
 p_{12}=\lambda^{-1},\quad p_{13}=\lambda,\quad p_{23}=\lambda^{-1}.
\]
Nakayama公式
$\mu(x_s)=(\prod_{a<s}p_{as})(\prod_{b>s}p_{sb}^{-1})x_s$
の三つの係数は全て1である．従ってcanonical bimoduleは$A_\lambda(-3)$であり，
\cite[Definition 1.3, Example 3.6]{MU}から$\qgr$上の反Serre関手は$(3)$である．

$d,c$は関係式を保ち，$d^3=c^3=1$である．次数シフトと可換するので三乗の等式が従う．
また$dc=\zeta cd$であるが，全次数1生成元を同じscalarで掛ける自己同型は，
次数$n$でscalarの$n$乗を掛ける自然同型によりgraded加群圏上恒等となる．
従って$d^*,c^*$は自然同型まで可換であり，三根も同様である．
$\tau$で割ると位数3の三つの可換元で生成されるので，位数は27を割る．
\end{proof}

\begin{proposition}[Hesse不変方向と半古典極限]\label{prop:quantum-tangent}
$q=1$での三根を同時に関手として持ち上げられる一次空間は
\[
 \HH^2(\mathbb P^2)^{s_0,s_d,s_c}
 =\operatorname{span}\{x^3+y^3+z^3,\ xyz\}.
\]
族\eqref{eq:quantum-family}の$q=1+t$方向の半古典括弧は
$\{x,y\}=xy$，$\{y,z\}=yz$，$\{z,x\}=zx$であり，
射影Poisson構造はこの2次元空間の$\C xyz$方向である．
従って上の族は，共通不変方向のうち少なくとも一つを明示的に代数化する．
\end{proposition}
\begin{proof}
$\HH^2(\mathbb P^2)=H^0(\OO(3))$でtensorの作用は恒等である．
命題\ref{prop:plane}により$d$の固定空間は$x^3,y^3,z^3,xyz$で張られる．
$c$は最初の三つを巡回し$xyz$を固定するので，共通固定部分は表示の2次元空間となる．
この一次主張には定理\ref{thm:fixed}を各関手へ適用する．
$\HH^1(\mathbb P^2)\ne0$なので定理\ref{thm:finite-package}は用いない．

括弧は関係式の$t$の一次係数から計算する．$x\ne0$上の$u=y/x,v=z/x$について
Leibniz則から$\{u,v\}=3uv$であり，反標準切断は非零scalar倍の$xyz$である．
三根の冪関係は全fiberで直接確認済みである．これはHesse不変空間全体の代数化でも，
$\HH^1\ne0$の場合の一般持上げ定理でもない．
\end{proof}

\begin{remark}[非可換な環と非可換な圏]
$\lambda\ne1$では$A_\lambda$は非可換だが，それだけで$\qgr A_\lambda$が
可換曲面の層圏と非同値であるとは結論しない．例えば$\lambda^3=1$では多項式環の
Zhang twistとなることがある．$q=1$での非零Poisson接方向，明示代数族の平坦性，
各fiber上の根の存在をそれぞれ別の主張として扱う．
\end{remark}

\section{安定性と成長率への二つの接続}
\label{sec:dynamics}
\begin{proposition}[Gepner位相の必要条件]\label{prop:gepner}
数値群$N(\mathcal D)_\Q$が有限次元で，$\tau_*$がunipotent，$s^r=\tau$とする．
数値的Bridgeland安定性$\sigma=(Z,\mathcal P)$が
$s\sigma=\sigma\cdot z$，$z\in\C$を満たすなら$rz\in2\Z$である．
さらに$s_*$自身がunipotentなら$z\in2\Z$である．
\end{proposition}
\begin{proof}
Toda\cite[Definition 2.3]{TodaGepner}の右$\C$作用の規約では
$Z\circ s_*^{-1}=e^{-i\pi z}Z$である．非零の$Z$は双対作用の固有ベクトルである．
$(s_*^{-1})^r=\tau_*^{-1}$の固有値は全て1なので$e^{-i\pi rz}=1$，
従って$rz\in2\Z$．$s_*$もunipotentなら既に$e^{-i\pi z}=1$となる．
\end{proof}

\begin{corollary}[Fano全圏におけるGepner型の非存在]\label{cor:no-gepner}
$X$を正次元の滑らかな射影複素Fano多様体とし，$s^r=(-)\otimes\omega_X^{-1}$とする．
数値群全体に関するsupport propertyを満たす数値的安定性$\sigma$で，
$s\sigma=\sigma\cdot z$，$z\in\C$となるものは存在しない．
\end{corollary}
\begin{proof}
線束tensorの数値作用はRiemann--Rochによりunipotentである．
前命題から$rz\in2\Z$なので，$\tau\sigma=\sigma\cdot rz$は全対象のmassを保つ．
support propertyをノルムの形で書くと，ある$C>0$について，全semistable対象$F$に
$\|[F]\|\le C|Z(F)|$が成立する．HN因子に分解し三角不等式を用いれば，全対象$E$で
$\|[E]\|\le C m_\sigma(E)$となる．従って
\[
 \|[\omega_X^{-m}]\|=\|[\tau^m\OO_X]\|
 \le C m_\sigma(\tau^m\OO_X)=C m_\sigma(\OO_X)
\]
は$m$によらず有界となる．一方，数値$K$群から数値因子群への第一Chern類を取ると
$c_1(\omega_X^{-m})=m(-K_X)$で，$-K_X$は正次元Fanoで非零だから非有界である．矛盾．
\end{proof}

\begin{remark}
系\ref{cor:no-gepner}は右$\C$作用に関するGepner型を扱う．より広い
$\widetilde{\mathrm{GL}}_2^+(\R)$作用によるSerre不変性を禁止した定理ではない．
またambientの全導来圏に関する結論であり，mutation後の残余圏やmatrix factorizationにはそのまま適用しない．
論証はsupport propertyと数値的成長の直接の帰結で，独立新規性は主張しない．
\end{remark}

\begin{remark}
これは存在定理ではなく必要条件である．正規化Fano根に非零の分数位相を持つGepner条件を
直接課すことはできない．matrix factorizationのgrade shift $\rho$に対する
$\rho^d=[2]$と位相$2/d$は別の関係である\cite{TodaGepner}．
Aravena--Liu--Tevelev\cite[Theorem 1.2]{ALT}の2026年11月のpreprintは
smooth hypersurfaceのGepner安定性の一般構成を主張する．本稿では主張を確認したが証明全体を監査しておらず，
上の命題にも以下の結果にもその存在定理は使わない．
\end{remark}

\begin{proposition}[Fano根のエントロピーは指数を区別しない]\label{prop:entropy}
滑らかな射影$n$次元複素Fano多様体上のheart-exactな根$s^r=(-)\otimes\omega_X^{-1}$について，
DHKKの圏論的エントロピーとFan--Fu--Ouchiの多項式エントロピーは，全$t\in\R$で
\[
 h_t(s)=0,\qquad h_t^{\mathrm{pol}}(s)=n
\]
である．正規化は仮定しない．
\end{proposition}
\begin{proof}
vector bundleであるsplit generator $G$を取る．$s=L\otimes f^*$なので$s^jG$もvector bundleである．
$k=rm+j$，$0\le j<r$と書くと
\[
 s^kG=(s^jG)\otimes\omega_X^{-m}.
\]
有限個の$j$にSerre消滅を適用すれば，$m\gg0$で全ての高次Extが消える．
Riemann--Rochにより
\[
 \dim\Hom(G,s^{rm+j}G)=
 \frac{\rk(G)^2(-K_X)^n}{n!}m^n+O(m^{n-1}),
\]
最高次係数は正で$j$によらない．従って全列のExt-distanceは正の定数倍の$k^n$に漸近する．
\cite[Theorem 2.6]{DHKK}および\cite[Lemma 2.7, \S6.2]{FFO}のExt-distance表示により両式を得る．
ここでは一般の多項式エントロピーの冪不変性を仮定せず，各剰余類で実際の極限を確認した．
\end{proof}

射影平面の二つの根は命題\ref{prop:entropy}で同じ$(0,2)$を返すが，
命題\ref{prop:plane}で変形不変部分の次元は異なる．数値指数・Hochschild指標・エントロピーを
単一の不変量へ混同しないことが，複数分野を結ぶ際に重要である．

\section{関連領域への接続と残る問題}
\label{sec:connections}
\paragraph{非可換射影幾何とNakayama自己同型．}
Reyes--Rogalski--Zhang\cite[Theorem 0.3]{RRZ}は左graded twistについて
$\mu_{A^{\widetilde\sigma}}=\mu_A\sigma^\ell\xi_{\mathrm{hdet}(\sigma)}^{-1}$を与える．
従ってAS指数$\ell$，grading shift，反Serre根を同一視する前にNakayama twistを処理する必要がある．
標準AS平面の根を調べる次の課題は，このnorm型関係と$\HH^2$作用を同時に計算することである．

\paragraph{高次表現論・Calabi--Yau completion・cluster圏．}
Haniharaのroot pairは逆双対化双加群の分数化を，生成threadと結び付ける\cite{Hanihara}．
本稿の固定部分はその変形方向の必要な候補を与えるが，root pairの半直交性や巡回不変性を代替しない．
数値根，関手根，生成付きroot pair，巡回不変なdg根を別の段階として扱う．

\paragraph{Lefschetz分解・残余圏・Torelli．}
Kuznetsov--Smirnovの$L_{\mathfrak a}s$は残余圏の偏極を作る\cite{KS}．
Kuznetsov--Perry\cite{KP}はcomplete intersectionのSerre関手にspherical twist補正が現れることを示す．
ambientの根の変形が，例外blockの変形と共に残余回転を持ち上げる条件を調べる必要がある．
LinのHochschild--Serre代数と，次数4のorbifold canonical ringによる復元\cite{Lin,LRZ}は，
単なる数値格子より多い情報が必要な理由を示す隣接理論である．

\paragraph{ミラー対称性・Poisson量子化・量子接続．}
Auroux--Katzarkov--Orlov\cite[Theorems 1.4, 1.6]{AKO}はdel Pezzoと非可換変形を
Lefschetz消滅サイクルおよび$B+i\omega$パラメータに結び付ける．
従って$\HH^2$の根の固定部分をmirror変形パラメータで表す問題は具体的である．
ただしmirror同値だけから，根をsymplectic monodromyの幾何学的分数と同定したり，
heart-exact性を保証したりはできない．
Karube--Robotis--Zuliani\cite{KRZ}は量子コホモロジーcentral chargeのliftと準収束経路を扱う．
本稿ではabstractを確認した段階であり，その安定性構成を使用していない．

\paragraph{モジュライ・marking忘却．}
次数4の配置は根という標識を残す変形と，圏だけを変形する場合の差を示す．
別研究\cite{Marking}のPicard再標識の障害と組み合わせ，根を持つ圏の相対群oid・慣性を記述することが次の課題である．
一つの根の固定部分を全自己同値群の商や大域モジュライと同一視しない．

\paragraph{t-構造付き形式moduliと原子的分解．}
Iwanari\cite[Theorems 4.5, 4.9, 4.12]{Iwanari}はcompleteな$t$-構造を保つ変形のHochschild制御を扱う．
一般のnaive Morita変形とflat Abel変形の差は重要であり，本稿はtilting比較の明示範囲に限定する．
またElagin--Schneider--Shinder\cite[Lemma 3.17]{ESS}は次数2の
$\mathcal B=\OO_X^\perp$について$\mathbb S_{\mathcal B}^2\simeq\gamma^*[3]$を記す．
これはambientの$L\otimes\gamma^*$の二乗が$\tau$となる関係とは異なる．
系\ref{cor:geiser-formal}の制約を，残余圏の全自己同値の非存在へ読み替えることはできない．
両文献はこの比較に用い，その一般定理を本稿の持上げ証明の入力にはしていない．

\paragraph{現在の証明境界．}
本稿はtilting bundleと$\HH^1=\HH^3=0$の下で，相対Serreと冪関係を含む全Artin次数の根の持上げを証明した．
Fanoの場合の相対反標準豊富性も得た．一方，$\HH^1\ne0$の一般の場合，有限反Serre商の仮定を外した複数根の同時持上げ，
root pairの巡回coherence，次数4・三次曲面の形式族の有限型代数化，および一般の非可換Fano指数上界は未解決である．関連文献の確認範囲・追加の接続案・検索上の不足は
\path{research/fano_index_connections.md}に整理した．
計算コード\path{computations/verify_root_deformations.py}と
\path{computations/verify_artin_roots.py}は指標，座標配置，affine平均化，非線型固定部分の厳密算術を検算する．
\path{computations/verify_root_packages.py}は$D_5$格子，全1024個の五頂点グラフ，
可換根の有限商，量子PBW積と共通指標を検算する．いずれも，
抽象的な圏・自己同値の存在を認証するものではない．



\begingroup\small
\begin{thebibliography}{99}
\bibitem{BO} A. Bondal and D. Orlov, \emph{Reconstruction of a variety from the derived category and groups of autoequivalences}, arXiv:alg-geom/9712029. Theorem 3.1.
\bibitem{KS} A. Kuznetsov and M. Smirnov, \emph{On residual categories for Grassmannians}, arXiv:1802.08097v3. Definition 2.1, Theorem 2.8.
\bibitem{Minamoto} H. Minamoto, \emph{Ampleness of two-sided tilting complexes}, author preprint, 7 March 2009; \url{https://www.math.kyoto-u.ac.jp/preprint/2009/11minamoto.pdf}. Definitions 0.1/2.6, Theorem 2.8.
\bibitem{MU} I. Mori and K. Ueyama, \emph{A categorical characterization of quantum projective spaces}, arXiv:1708.00167v4 (2021). Theorem 4.1. 旧v1の訂正前主張は使わない．
\bibitem{LO} K. Loginov and D. Orlov, \emph{Non-commutative Fano varieties and three-dimensional minifolds}, arXiv:2609.39766 (2026). Definitions 2.1, 2.6; Proposition 2.7.
\bibitem{Hanihara} N. Hanihara, \emph{Calabi--Yau completions for roots of dualizing dg bimodules}, arXiv:2412.18753v1 (2024). Definition 2.5, Definitions 3.1--3.2, Theorem 3.3.
\bibitem{Perry} A. Perry, \emph{Hochschild cohomology and group actions}, arXiv:1807.09268v2. Proposition 1.2 and \S5.
\bibitem{Toda} Y. Toda, \emph{Deformations and Fourier--Mukai transforms}, J. Differential Geom. \textbf{81} (2009), 197--224; arXiv:math/0502571v3. Theorem 1.1.
\bibitem{LVdB} W. Lowen and M. Van den Bergh, \emph{Hochschild cohomology of abelian categories and ringed spaces}, arXiv:math/0405227. Theorem 3.1 and \S4.1.
\bibitem{HP} L. Hille and M. Perling, \emph{Tilting Bundles on Rational Surfaces and Quasi-Hereditary Algebras}, Ann. Inst. Fourier \textbf{64} (2014), 625--644; arXiv:1110.5843. Theorem 1.1.
\bibitem{BW} S. Barmeier and Z. Wang, \emph{Deformations of categories of coherent sheaves via quivers with relations}, Algebraic Geometry \textbf{11} (2024), 1--36; DOI:10.14231/AG-2024-001; arXiv:2107.07490. Theorems 3.1, 3.18, Proposition 3.21.
\bibitem{LowenOb} W. Lowen, \emph{Obstruction theory for objects in abelian and derived categories}, Comm. Algebra \textbf{33} (2005), 3195--3223; arXiv:math/0407019. Theorems A, B; arXiv版Theorem 6.10．本文ではflat/Tor版を用いる．
\bibitem{KellerDIH} B. Keller, \emph{Derived invariance of higher structures on the Hochschild complex}, preprint, August 2003, last modified 16 May 2018. \url{https://webusers.imj-prg.fr/~bernhard.keller/publ/dih.pdf}. Theorem 3.2, \S4.6.
\bibitem{BDW} H. Bursztyn, V. Dolgushev and S. Waldmann, \emph{Morita equivalence and characteristic classes of star products}, arXiv:0909.4259. Proposition 3.7, Theorem 3.11.
\bibitem{DM} I. Dolgachev and G. Martin, \emph{Automorphisms of del Pezzo surfaces in odd characteristic}, arXiv:2303.16170v2; J. Lond. Math. Soc. \textbf{109} (2024), e12905. \S2.3, \S4.
\bibitem{HKT} B. Hassett, A. Kresch and Y. Tschinkel, \emph{On the moduli of degree 4 Del Pezzo surfaces}, Adv. Stud. Pure Math. \textbf{69} (2016), 349--386. Author text: \url{https://cims.nyu.edu/~tschinke/papers/yuri/13noheighttwo/noheighttwo13.pdf}. Introduction, \S2.
\bibitem{TodaGepner} Y. Toda, \emph{Gepner type stability conditions on graded matrix factorizations}, arXiv:1302.6293. Definition 2.3, \S2.2.
\bibitem{ALT} A. Aravena, P. Liu and J. Tevelev, \emph{Gepner stability conditions on hypersurface categories}, arXiv:2611.05350 (2026). Theorem 1.2の主張のみを比較に使用；証明全体未監査．
\bibitem{DHKK} G. Dimitrov, F. Haiden, L. Katzarkov and M. Kontsevich, \emph{Dynamical systems and categories}, arXiv:1307.8418v1. Theorem 2.6.
\bibitem{FFO} Y.-W. Fan, L. Fu and G. Ouchi, \emph{Categorical polynomial entropy}, Adv. Math. \textbf{383} (2021), 107655. Lemmas 2.7, 2.10; Proposition 6.4, Theorem 6.7.
\bibitem{RRZ} M. Reyes, D. Rogalski and J. J. Zhang, \emph{Skew Calabi--Yau algebras and homological identities}, arXiv:1302.0437v2. Theorem 0.3.
\bibitem{KP} A. Kuznetsov and A. Perry, \emph{Serre functors and dimensions of residual categories}, arXiv:2109.02026. Introduction, \S6.
\bibitem{Lin} X. Lin, \emph{Group action of Hochschild--Serre algebra and categorical reconstruction}, arXiv:2609.07562 (2026). \S2--3の定義・主張を比較．
\bibitem{LRZ} X. Lin, M. Rampazzo and S. Zhang, \emph{Categorical reconstruction of del Pezzo surfaces: Hochschild--Serre algebras and spinor modifications}, arXiv:2609.10344 (2026). Introduction, \S4の主張を比較．
\bibitem{AKO} D. Auroux, L. Katzarkov and D. Orlov, \emph{Mirror symmetry for Del Pezzo surfaces: Vanishing cycles and coherent sheaves}, arXiv:math/0506166. Theorems 1.4, 1.6.
\bibitem{KRZ} T. Karube, A.-A. Robotis and V. Zuliani, \emph{Toward the noncommutative minimal model program for Fano varieties}, arXiv:2601.20739 (2026). 本稿ではabstractのみ確認．
\bibitem{Iwanari} I. Iwanari, \emph{Hochschild cohomology and deformation theory of stable infinity-categories with t-structures}, arXiv:2506.21867v1 (2025). Introduction, Theorems 4.5, 4.9, 4.12の主張を比較；全証明未監査．
\bibitem{ESS} A. Elagin, J. Schneider and E. Shinder, \emph{Atomic decompositions for derived categories of $G$-surfaces}, arXiv:2512.05064v1 (2025). Lemma 3.17とその短い帰着証明．
\bibitem{ATVSk} M. Artin, J. Tate and M. Van den Bergh,
\emph{Some algebras associated to automorphisms of elliptic curves},
The Grothendieck Festschrift, Vol. I, Progr. Math. \textbf{86} (1990), 33--85.
\bibitem{MoriSmith} I. Mori and S. P. Smith,
\emph{The classification of 3-Calabi--Yau algebras with 3 generators and 3 quadratic relations},
arXiv:1502.07403. Cyclic potentials and regularity.
\bibitem{DeLaet} K. De Laet,
\emph{On the center of 3-dimensional and 4-dimensional Sklyanin algebras},
J. Algebra \textbf{487} (2017), 244--268.
\url{https://arxiv.org/abs/1612.06158}. Heisenberg symmetry.
\bibitem{Index} 本研究の現行原稿，\emph{正規化された圏論的Fano指数}, \texttt{papers/normalized\_fano\_index.tex}, v14．
\bibitem{Marking} 関連研究の現行原稿，\emph{非可換三次曲面のマーキング忘却とPicard作用}, \texttt{uedakazushi/nc-cubic-ambient-moduli}, surface-v6，2026年11月17日．一次Picard障害との接続のみ参照し，本稿の証明の外部入力にはしない．
\end{thebibliography}
\endgroup
\end{document}
